\documentclass{article}
\usepackage{graphicx} % Required for inserting images

\title{Homework for the course of Methodology and Research Methods of Political Science (Questions 1-3)}
\author{Vsevolod Muronets}
\date{The 5th of December 2023}

\begin{document}

\maketitle 
\textbf{Problem  1.}
\\

1.	 First syllogism is a \textbf{denial of antecedent}, so the argument is logically \textbf{invalid}. 

2.	Second syllogism is a \textbf{denial of consequent}, and argument is logically \textbf{valid}.

3.	 Third – \textbf{affirming the antecedent}, which makes a logically \textbf{valid} argument.

4.	 In the fourth syllogism there is an \textbf{affirmation of the consequent}, so the argument is \textbf{invalid}. 
\\

\textbf{Problem 2.}
\\

1.	A is \textbf{sufficient} for B, since if x is more than 50 we always observe that it is more than 45. But it is \textbf{not necessary} since we can observe x being more than 45 and less than 50, like when it is equal to 49, for example. 

2.	A is \textbf{necessary} for B, because it is impossible to cast an invalid vote without being voting. But it is \textbf{not sufficient} because one can turn out to vote and cast a valid vote. 

3.	A is \textbf{sufficient} for B, because having a cold is enough to be sick, but it is \textbf{not necessary}, because one can be sick for other reasons.

4.	 A is \textbf{sufficient} for B, because when there is a civil war there is always violent conflict, but the latter can happen for some other reasons, so it is \textbf{not necessary}. 

5.	A is \textbf{necessary} for B, since if we consider the distinction between autocracy and democracy, there can be no democracy with less than two parties, but it is \textbf{not sufficient} because in a autocracy there can be, for example, one ruler without parties, or an autocratic ruler can make a various-party system but as a mere façade without any real power.

6.	A is \textbf{both sufficient and necessary} for B, since when there is a large number of veto players in a country there is always high level of policy stability (sufficient), and there can be no high levels of policy stability when there are only a few veto players who tend to make policy in their own interest (necessary).

7.	A is \textbf{sufficient} for B, because a protest is a type of collective action, but not necessary,  since some other types of collective actions may occur without any relation to protests.  
\\

\textbf{Problem 3.}
\\

The \textbf{research question} (puzzle) I would like to study, besides my ongoing research, is a question of why some street protest movements achieve at least some of their stated goals and some never do (in a shorter form, why some street protest work and others do not). \textbf{My theory} would be that actual participation in street movement compared to simple verbal, “clicktivistical” or any other formal approval of the movement, combined with sustainable organization, contributes to the success of street protests. From this theory I can derive \textbf{two testable hypotheses}: 1) \textbf{that when the relation between} the number of people who actually appear “on the streets”, that is, directly participate in the street protest, and the amount of users that support the protest on social media, \textbf{is high}, the movement have higher possibility to succeed; 2) and that those street movements that end up forming some organization (like a party, or a non-profit organization, a sustainable community etc.) are more successful than those which end with the end of the actual street protest, even though the latter may be very popular among people. 


\end{document}